% !TEX program = xelatex
% 컴파일 방법: xelatex pde_notes.tex (두 번 실행하면 목차가 갱신됩니다)
%
% 필요 패키지: fontspec (XeLaTeX 기본), amsmath, amssymb, geometry, hyperref, xcolor, tcolorbox, booktabs
% TeX Live 2020 이상 권장. Noto Sans CJK KR 폰트가 시스템에 설치되어 있어야 합니다.
% kotex가 있는 환경이라면 아래의 fontspec 설정을 \usepackage{kotex}으로 대체해도 됩니다.

\documentclass[11pt,a4paper]{article}

\usepackage{fontspec}
\setmainfont{Noto Sans CJK KR}[Scale=0.95]
\setsansfont{Noto Sans CJK KR}[Scale=0.95]
\setmonofont{Noto Sans Mono CJK KR}[Scale=0.9]

\usepackage{amsmath, amssymb, amsthm}
\usepackage[margin=2.5cm]{geometry}
\usepackage{hyperref}
\usepackage{xcolor}
\usepackage[most]{tcolorbox}
\usepackage{booktabs}
\usepackage{array}
\usepackage{enumitem}
\usepackage{parskip}

\hypersetup{
  colorlinks=true,
  linkcolor=blue!60!black,
  urlcolor=blue!60!black,
  citecolor=blue!60!black,
}

\newtcolorbox{insightbox}[1][]{
  colback=blue!5!white,
  colframe=blue!50!black,
  fonttitle=\bfseries,
  title=#1,
  boxrule=0.5pt,
}

\newtcolorbox{warningbox}[1][]{
  colback=orange!5!white,
  colframe=orange!50!black,
  fonttitle=\bfseries,
  title=#1,
  boxrule=0.5pt,
}

\newtcolorbox{historybox}[1][]{
  colback=gray!5!white,
  colframe=gray!50!black,
  fonttitle=\bfseries,
  title=#1,
  boxrule=0.5pt,
}

\title{\textbf{편미분방정식 풀이의 280년}\\[0.3em] \large 달랑베르에서 FNO까지 --- 수리물리학과 계산과학의 흐름}
\author{대화 학습 노트}
\date{2026년 5월}

\begin{document}

\maketitle

\begin{abstract}
이 문서는 편미분방정식(PDE)의 기본 풀이 전략에서 시작하여, 19세기의 적분변환과 그린 함수, 20세기의 수치해석, 그리고 21세기의 Physics-Informed Neural Networks (PINN)와 Fourier Neural Operator (FNO)에 이르기까지 280년에 걸친 수리물리학과 계산과학의 흐름을 정리한 학습 노트이다. 1747년 달랑베르의 파동방정식 풀이에서 출발하여, 라플라스ㆍ푸리에의 적분변환, 그린의 점원천 응답 이론, 수치해석의 차분 근사, 그리고 신경연산자에 의한 함수 공간 학습까지---모든 도구는 결국 ``다루기 어려운 미분 연산을 다루기 쉬운 무언가로 바꿔치기한다''는 단 하나의 사상에 수렴한다.
\end{abstract}

\tableofcontents
\newpage

% ============================================================
\section{서론: 왜 편미분방정식인가}
% ============================================================

자연계의 거의 모든 동역학은 편미분방정식으로 기술된다. 파동의 전파, 열의 확산, 유체의 흐름, 전자기장의 변화, 양자역학의 슈뢰딩거 방정식, 일반상대성이론의 아인슈타인 방정식, 금융수학의 블랙-숄즈 방정식---모두 PDE이다. 따라서 ``PDE를 푼다''는 것은 곧 ``자연을 예측한다''는 것과 거의 동의어이다.

이 문서는 PDE 풀이의 기본 개념에서 출발하여 시간 순서로 도구의 발전사를 따라간다. 각 장은 독립적으로 읽힐 수 있지만, 전체를 순서대로 따라가면 ``미분이라는 까다로운 연산을 다루기 쉬운 형태로 변환하는 기술''이 어떻게 발전해 왔는지가 한 줄기로 이어진다.

% ============================================================
\section{편미분방정식 풀이의 개요}
% ============================================================

\subsection{ODE와 PDE의 결정적 차이}

상미분방정식(ODE)과 편미분방정식(PDE)의 가장 중요한 차이는 \emph{해의 자유도}에 있다.

\begin{itemize}
\item ODE: $\dfrac{du}{dx} = 0 \Rightarrow u = C$ \quad (상수 하나가 자유)
\item PDE: $\dfrac{\partial u}{\partial x} = 0 \Rightarrow u = f(y)$ \quad (\textbf{함수 하나가 통째로 자유})
\end{itemize}

PDE의 해 공간이 무한히 더 크기 때문에, 특정 해를 결정하려면 \textbf{경계조건} 또는 \textbf{초기조건}이 반드시 필요하다. 이것이 PDE 풀이를 ODE보다 본질적으로 까다롭게 만드는 첫 번째 요인이다.

\subsection{대표적 풀이 기법}

\begin{itemize}
\item \textbf{변수분리법(Separation of Variables)}: $u(x,t) = X(x)T(t)$로 가정하여 PDE를 여러 ODE로 분해. 열방정식, 라플라스 방정식에서 자주 사용.
\item \textbf{특성곡선법(Method of Characteristics)}: 해가 일정하게 유지되는 특별한 곡선을 따라가며 PDE를 ODE로 환원. 1차 PDE와 쌍곡형 PDE에 효과적.
\item \textbf{좌표변환}: 새로운 변수를 도입해 방정식을 단순화. 달랑베르의 파동방정식 풀이가 대표적.
\item \textbf{적분변환}: 푸리에ㆍ라플라스 변환으로 미분을 곱셈으로 변환.
\item \textbf{그린 함수}: 점원천의 응답을 미리 구하고 적분으로 일반 해를 얻음.
\item \textbf{수치해석}: 컴퓨터로 근사해를 얻음 (FDM, FVM, FEM 등).
\end{itemize}

\subsection{PDE의 분류}

2계 선형 PDE는 판별식 $B^2 - AC$의 부호로 세 종류로 분류된다.

\begin{center}
\begin{tabular}{lll}
\toprule
유형 & 조건 & 대표 예 \\
\midrule
타원형(elliptic) & $B^2 - AC < 0$ & 라플라스 방정식 $\nabla^2 u = 0$ \\
포물형(parabolic) & $B^2 - AC = 0$ & 열방정식 $u_t = \alpha u_{xx}$ \\
쌍곡형(hyperbolic) & $B^2 - AC > 0$ & 파동방정식 $u_{tt} = c^2 u_{xx}$ \\
\bottomrule
\end{tabular}
\end{center}

% ============================================================
\section{달랑베르의 파동방정식과 연쇄법칙}
% ============================================================

\subsection{역사적 배경}

1747년, 장 르 롱 달랑베르(Jean le Rond d'Alembert)는 양 끝이 고정된 진동하는 줄(예: 바이올린 현)의 운동을 분석하면서 다음 편미분방정식에 도달했다.

\begin{equation}
\frac{\partial^2 u}{\partial t^2} = c^2 \frac{\partial^2 u}{\partial x^2}
\label{eq:wave}
\end{equation}

여기서 $u(x,t)$는 위치 $x$, 시간 $t$에서 줄의 변위, $c$는 파동의 전파 속도이다. 이것이 1차원 파동방정식이며, 수리물리학에서 편미분방정식을 본격적으로 사용한 거의 최초의 사례로 꼽힌다.

\subsection{달랑베르의 풀이}

달랑베르는 다음과 같은 새 좌표를 도입했다.

\begin{equation}
\xi = x - ct, \qquad \eta = x + ct
\end{equation}

\subsubsection{1단계: 연쇄법칙으로 미분 변환}

다변수 연쇄법칙은 다음과 같다. $u = u(\xi, \eta)$, $\xi = \xi(x,t)$, $\eta = \eta(x,t)$일 때,

\begin{equation}
\frac{\partial u}{\partial x} = \frac{\partial u}{\partial \xi}\frac{\partial \xi}{\partial x} + \frac{\partial u}{\partial \eta}\frac{\partial \eta}{\partial x}
\end{equation}

직관적으로 ``모든 의존 경로의 합''으로 이해할 수 있다. $x$가 살짝 움직이면 $\xi$도 움직이고, $\eta$도 움직이며, 이 둘이 각각 $u$를 움직인다.

새 좌표의 편도함수는 다음과 같이 계산된다.

\begin{equation}
\frac{\partial \xi}{\partial x} = 1, \quad \frac{\partial \xi}{\partial t} = -c, \quad \frac{\partial \eta}{\partial x} = 1, \quad \frac{\partial \eta}{\partial t} = c
\end{equation}

따라서 미분 연산자는 다음과 같이 변환된다.

\begin{equation}
\frac{\partial}{\partial x} = \frac{\partial}{\partial \xi} + \frac{\partial}{\partial \eta}, \qquad \frac{\partial}{\partial t} = -c\frac{\partial}{\partial \xi} + c\frac{\partial}{\partial \eta}
\end{equation}

\subsubsection{2단계: 2계 미분 계산}

연산자를 두 번 적용하면, 마치 $(a+b)^2$, $(-a+b)^2$를 전개하듯 다뤄진다.

\begin{align}
\frac{\partial^2 u}{\partial x^2} &= u_{\xi\xi} + 2u_{\xi\eta} + u_{\eta\eta} \\
\frac{\partial^2 u}{\partial t^2} &= c^2\bigl(u_{\xi\xi} - 2u_{\xi\eta} + u_{\eta\eta}\bigr)
\end{align}

여기서 혼합편미분의 순서 무관성(\textbf{클레로 정리}): $u_{\xi\eta} = u_{\eta\xi}$ 를 사용했다.

\subsubsection{3단계: 원래 식에 대입}

식 (\ref{eq:wave})에 대입하면

\begin{equation}
c^2(u_{\xi\xi} - 2u_{\xi\eta} + u_{\eta\eta}) = c^2(u_{\xi\xi} + 2u_{\xi\eta} + u_{\eta\eta})
\end{equation}

같은 항을 소거하면

\begin{equation}
\frac{\partial^2 u}{\partial \xi \partial \eta} = 0
\end{equation}

\subsubsection{4단계: 적분과 일반해}

$\eta$에 대해 적분하면 $\partial u/\partial \xi$는 $\xi$만의 함수, 즉 $F(\xi)$. 다시 $\xi$로 적분하면 적분상수는 $\eta$만의 함수가 되므로

\begin{equation}
u(\xi, \eta) = f(\xi) + g(\eta)
\end{equation}

원래 변수로 복귀하면 \textbf{달랑베르의 일반해}:

\begin{empheq}[box=\fbox]{equation}
u(x, t) = f(x - ct) + g(x + ct)
\end{empheq}

\subsection{물리적 해석}

\begin{itemize}
\item $f(x - ct)$: 모양을 그대로 유지한 채 속도 $c$로 \textbf{오른쪽으로 진행}하는 파동
\item $g(x + ct)$: 같은 속도로 \textbf{왼쪽으로 진행}하는 파동
\end{itemize}

시간이 $\Delta t$ 흐르면 $f$의 그래프 전체가 $c\Delta t$만큼 오른쪽으로 평행이동할 뿐 형태는 변하지 않는다.

\subsection{달랑베르 공식}

초기조건 $u(x, 0) = \varphi(x)$, $u_t(x, 0) = \psi(x)$가 주어지면 $f, g$를 명시적으로 풀 수 있다.

\begin{equation}
u(x, t) = \frac{\varphi(x-ct) + \varphi(x+ct)}{2} + \frac{1}{2c}\int_{x-ct}^{x+ct} \psi(s)\,ds
\label{eq:dalembert}
\end{equation}

수치해석도 푸리에 급수도 없이, 단지 좌표를 잘 잡아서 방정식의 본래 구조(특성곡선 $x \pm ct = \text{const}$)를 드러낸 것만으로 일반해가 곧장 나왔다.

% ============================================================
\section{적분변환: 라플라스와 푸리에}
% ============================================================

\subsection{핵심 아이디어: 어려운 세계에서 쉬운 세계로}

적분변환의 전체 전략은 다음과 같다.

\begin{center}
미분방정식 $\xrightarrow{\text{변환}}$ 대수방정식 $\xrightarrow{\text{풀이}}$ 해(변환된 형태) $\xrightarrow{\text{역변환}}$ 진짜 해
\end{center}

큰 수의 곱셈이 어려우니 로그를 취해 덧셈으로 바꿔 계산하고 다시 지수를 취해 돌아오는 것과 같은 발상이다.

\subsection{라플라스 변환의 정의}

함수 $f(t)$를 받아 새로운 함수 $F(s)$를 만드는 변환:

\begin{equation}
F(s) = \mathcal{L}[f](s) = \int_0^\infty e^{-st} f(t)\, dt
\label{eq:laplace}
\end{equation}

\subsection{핵심 성질: 미분이 곱셈으로}

\begin{equation}
\mathcal{L}\!\left[\frac{df}{dt}\right] = sF(s) - f(0)
\end{equation}

증명은 부분적분 한 줄이다.

\begin{equation}
\int_0^\infty e^{-st} f'(t)\, dt = \underbrace{[e^{-st}f(t)]_0^\infty}_{= -f(0)} + s\int_0^\infty e^{-st} f(t)\, dt = sF(s) - f(0)
\end{equation}

\subsection{간단한 ODE 시연}

다음을 풀어보자.

\begin{equation}
\frac{dy}{dt} + 2y = 0, \qquad y(0) = 3
\end{equation}

\textbf{1단계}: 양변에 라플라스 변환

\begin{equation}
(sY(s) - 3) + 2Y(s) = 0
\end{equation}

\textbf{2단계}: 대수적 풀이

\begin{equation}
Y(s) = \frac{3}{s+2}
\end{equation}

\textbf{3단계}: 역변환 (변환표: $1/(s+a) \to e^{-at}$)

\begin{equation}
y(t) = 3e^{-2t}
\end{equation}

미분방정식이 변환 → 대수 → 역변환의 세 단계로 풀렸다.

\subsection{파동방정식에 적용}

식 (\ref{eq:wave})에 $t$에 대해서만 라플라스 변환을 적용한다.

\begin{equation}
U(x, s) = \int_0^\infty e^{-st} u(x, t)\, dt
\end{equation}

$t$ 미분은 $s$ 곱셈이 되고, $x$ 미분은 그대로 살아남는다.

\begin{align}
\mathcal{L}\!\left[\frac{\partial^2 u}{\partial t^2}\right] &= s^2 U(x, s) - s\varphi(x) - \psi(x) \\
\mathcal{L}\!\left[\frac{\partial^2 u}{\partial x^2}\right] &= \frac{\partial^2 U}{\partial x^2}
\end{align}

원래 방정식은 $x$만의 ODE로 변환된다.

\begin{equation}
c^2 \frac{d^2 U}{dx^2} - s^2 U = -s\varphi(x) - \psi(x)
\end{equation}

\textbf{PDE가 차원이 하나 줄어 ODE가 됐다.} 이를 풀고 역변환하면 결국 달랑베르 공식 (\ref{eq:dalembert})와 같은 답이 나온다. 길은 달라도 종착지는 같다.

% ============================================================
\section{라플라스의 생애와 변환의 역사}
% ============================================================

\subsection{라플라스 변환은 라플라스가 만들었나}

이름은 라플라스의 것이지만 실제 역사는 더 복잡하다.

\begin{historybox}[역사의 흐름]
\begin{itemize}
\item \textbf{1744}: 오일러가 이런 형태의 적분을 미분방정식의 해로서 처음 연구
\item \textbf{1782}: 라플라스가 오일러의 정신을 이어받아 확률론 연구에 사용
\item \textbf{1809}: 라플라스가 변환을 확장하여 무한 확산 해를 다룸
\item \textbf{1821}: 코시가 라플라스 변환의 연산자 미적분을 개발
\item \textbf{1880년대}: 영국의 자가학습자 올리버 헤비사이드가 독자적으로 같은 도구에 도달, 전기회로에 응용
\item \textbf{1910}: 베이트먼이 현대적 형태의 라플라스 변환 사용
\item \textbf{1920년대}: 도이치가 체계화
\item \textbf{1937}: 처음으로 교과서 주제가 됨
\item \textbf{1950년경}: 헤비사이드의 연산자 미적분을 완전히 대체
\end{itemize}
\end{historybox}

라플라스가 변환을 \emph{발명}한 것은 아닐지 몰라도, 선배들의 산발적 아이디어를 훨씬 뛰어넘는 체계적 이론으로 만들어낸 공로는 분명히 그의 것이다. 우리가 아는 모습이 된 것은 사실상 \textbf{제2차 세계대전 후}이다.

\subsection{왜 이런 변환을 만들었는가}

\begin{itemize}
\item \textbf{확률론의 모함수(generating function)}: 라플라스는 동전 던지기, 도박 같은 확률 문제에서 이산 수열 $a_n$을 $\sum a_n x^n$ 같은 생성함수로 묶어 미적분으로 처리했다. 적분 형태의 변환은 이걸 연속 분포로 확장한 자연스러운 산물이었다.
\item \textbf{천체역학과 확산 문제}: 라플라스의 평생 과업은 태양계의 안정성 증명이었다. 뉴턴이 ``신의 주기적 개입이 필요하다''고까지 말한 문제를 순수 수학으로 반박하려 했다. 1809년에 자신의 변환을 확장하여 공간 전체로 무한히 확산되는 해를 찾았다.
\end{itemize}

\subsection{라플라스의 생애 요약}

\begin{itemize}
\item \textbf{1749년 3월 23일}: 노르망디 보몽-앙-오주 출생. 농부의 아들.
\item \textbf{1766}: 캉 대학 입학, 사제가 되기 위한 신학 공부
\item \textbf{1768}: 학위 없이 파리로. 달랑베르의 인정을 받아 에콜 밀리테르 수학 교수
\item \textbf{1773}: 24세에 프랑스 과학원 회원
\item \textbf{1780}: 라부아지에와 함께 호흡이 연소의 일종임을 증명
\item \textbf{1785}: 16세의 나폴레옹 보나파르트를 시험관 자격으로 만남 --- 운명적 인연
\item \textbf{1793}: 공포정치 시기 파리 밖으로 피신, 처형 모면 (라부아지에는 1794년 단두대)
\item \textbf{1799}: 나폴레옹에 의해 내무장관 임명 --- 6주 만에 해임
\item \textbf{1799--1825}: \emph{천체역학(Mécanique céleste)} 5권 출간
\item \textbf{1812}: \emph{확률에 관한 해석적 이론} 출간
\item \textbf{1814}: 부르봉 왕가에 충성을 갈아탐, 후작 작위 받음
\item \textbf{1827년 3월 5일}: 파리에서 사망. 과학 아카데미가 그날의 회의를 취소하여 경의 표함.
\end{itemize}

\subsection{유명한 일화: ``신은 필요 없는 가설입니다''}

라플라스가 \emph{천체역학}을 나폴레옹에게 헌정했을 때, 황제는 ``책에 신의 이름이 한 번도 나오지 않는다''고 지적했다. 라플라스의 답:

\begin{quote}
\itshape ``폐하, 그 가설은 필요하지 않았습니다 (Je n'ai pas eu besoin de cette hypothèse).''
\end{quote}

이 대화의 정확한 정황은 후대에 윤색되었다는 설이 강하지만, 발언의 정신은 라플라스의 결정론적 세계관을 잘 보여준다. 뉴턴이 신의 주기적 개입을 필요로 봤던 곳에서, 라플라스는 순수한 수학으로 태양계의 안정성을 증명해 보였다.

나폴레옹은 라플라스의 행정관 자격을 평하며 ``\textbf{라플라스는 미적분의 정신을 행정에 들고 왔다}''라고 했다 --- 모든 사안에 무한소 분석을 적용하려 했다는 비꼼이었다.

% ============================================================
\section{지수함수의 본질: 왜 e의 적분인가}
% ============================================================

\subsection{핵심 한 줄: $e^{ax}$는 미분의 고유함수}

\begin{equation}
\frac{d}{dx}e^{ax} = a \cdot e^{ax}
\end{equation}

즉 $e^{ax}$를 미분하면 형태가 변하지 않고 상수 $a$가 곱해질 뿐이다. \textbf{미분이 곱셈이 된다.}

선형대수의 언어로, 미분 연산자 $d/dx$를 행렬로 보면 $e^{ax}$는 그 행렬의 \textbf{고유벡터(eigenvector)}이고 $a$는 \textbf{고유값(eigenvalue)}이다. 행렬을 대각화하면 모든 계산이 쉬워지듯, 미분을 다룰 때 $e^{ax}$를 기저로 삼으면 모든 게 쉬워진다.

이게 모든 적분변환에서 $e$가 등장하는 \textbf{근본 이유}이다. 다른 함수로는 이게 안 된다.

\begin{itemize}
\item $\frac{d}{dx}\sin x = \cos x$ --- 모양이 바뀐다
\item $\frac{d}{dx}x^n = nx^{n-1}$ --- 모양과 차수가 바뀐다
\item $\frac{d}{dx}e^{ax} = a\,e^{ax}$ --- \textbf{모양 그대로, 상수만 곱해진다}
\end{itemize}

\subsection{왜 곱해서 적분하는가: 내적의 개념}

함수의 내적은 벡터 내적의 자연스러운 확장이다.

\begin{equation}
\langle f, g \rangle = \int f(t)\, g(t)\, dt
\end{equation}

이산적인 합($\sum$)이 연속적인 합($\int$)으로 바뀐 것뿐, 본질은 같다. \textbf{함수도 무한차원 공간의 벡터로 볼 수 있다}는 게 20세기 수학의 핵심 관점이다.

\begin{itemize}
\item 푸리에 변환: $\hat{f}(\omega) = \int f(t)\, e^{-i\omega t}\, dt$ --- ``$f$가 주파수 $\omega$의 진동 방향에 얼마나 들어있는가''
\item 라플라스 변환: $F(s) = \int_0^\infty f(t)\, e^{-st}\, dt$ --- ``$f$가 감쇠율 $s$의 지수 방향에 얼마나 들어있는가''
\end{itemize}

적분변환은 \textbf{함수를 지수함수라는 기저로 분해}하는 작업이다. 빛을 프리즘에 통과시켜 무지개로 분해하듯, 함수를 무수히 많은 $e^{\text{something}}$ 성분으로 펼치는 것.

\subsection{푸리에와 라플라스의 통합 관점}

오일러 공식 $e^{i\theta} = \cos\theta + i\sin\theta$가 의미하는 바: \textbf{사인ㆍ코사인은 사실 $e^{i\theta}$의 변형}이다.

\begin{equation}
\cos\theta = \frac{e^{i\theta} + e^{-i\theta}}{2}, \qquad \sin\theta = \frac{e^{i\theta} - e^{-i\theta}}{2i}
\end{equation}

따라서 통합된 관점:

\begin{insightbox}[푸리에와 라플라스의 통합]
\textbf{푸리에든 라플라스든, 결국 함수를 $e^{ax}$ 형태의 지수함수들로 분해한다. 다만 어떤 $a$를 허용하느냐가 다르다.}

\begin{center}
\begin{tabular}{lll}
\toprule
변환 & ``$a$''의 범위 & 의미 \\
\midrule
푸리에 & 순허수 ($a = i\omega$) & 순수 진동 (감쇠 없음) \\
라플라스 & 복소수 ($a = -s$, $s = \sigma + i\omega$) & 진동 + 감쇠 \\
\bottomrule
\end{tabular}
\end{center}

라플라스 변환은 푸리에 변환의 확장이다.
\end{insightbox}

\subsection{한 개의 지수함수 vs 무한합}

\textbf{$y = x^2$ = 한 개의 $e^{ax}$}? --- 불가능하다. 그래프 모양이 완전히 다르다.

\textbf{$y = x^2$ = 무한히 많은 $e^{ax}$의 가중합}? --- 적절한 조건 하에서 가능하다.

예: $[-\pi, \pi]$ 구간에서

\begin{equation}
x^2 = \frac{\pi^2}{3} + 4\sum_{n=1}^{\infty} \frac{(-1)^n}{n^2}\cos(nx)
\end{equation}

이 공식은 실제로 오일러가 유명한 $\sum 1/n^2 = \pi^2/6$을 증명할 때 쓴 도구이기도 하다.

\subsection{분해 가능성의 경계}

함수의 종류와 적절한 변환 도구의 대응:

\begin{center}
\begin{tabular}{p{6cm}p{8cm}}
\toprule
함수의 종류 & 적절한 변환 \\
\midrule
적당히 부드럽고 무한대에서 0으로 가는 함수 & 푸리에 변환 \\
한쪽 방향으로만 정의되고 $e^{-st}$가 감쇠시킬 정도로 발산 & 라플라스 변환 \\
구간 위에서 정의된 함수 & 푸리에 급수 \\
상수, 다항식, 디랙 델타 등 일반 함수가 아닌 것 & 분포(distribution) 이론 \\
\bottomrule
\end{tabular}
\end{center}

푸리에가 ``\emph{모든} 함수가 사인파의 합''이라고 외쳤을 때 라그랑주가 격분한 것은 부분적으로 옳았다. 진짜 \emph{모든} 함수는 아니다. 하지만 물리적으로 의미 있는 함수들의 매우 넓은 집합은 정말로 분해된다.

% ============================================================
\section{부분적분과 통찰력 있는 질문들}
% ============================================================

\subsection{$\mathcal{L}[t^2]$의 계산}

목표:

\begin{equation}
\mathcal{L}[t^2] = \int_0^{\infty} t^2 e^{-st}\, dt = \frac{2}{s^3}
\end{equation}

부분적분 공식 $\int u\, dv = uv - \int v\, du$를 적용한다. 다항식은 미분하면 단순해지므로 $u = t^2$로 잡는다.

\subsubsection{1단계 부분적분}

$u = t^2,\ du = 2t\,dt,\ dv = e^{-st}\,dt,\ v = -\frac{1}{s}e^{-st}$

\begin{equation}
\int_0^\infty t^2 e^{-st}\, dt = \underbrace{\left[-\frac{t^2}{s}e^{-st}\right]_0^\infty}_{= 0} + \frac{2}{s}\int_0^\infty t\, e^{-st}\, dt
\end{equation}

경계항이 0인 이유: $t \to \infty$일 때 $e^{-st}$가 다항식 $t^2$보다 훨씬 빠르게 0으로 가기 때문 ($s > 0$ 가정).

\subsubsection{2단계 부분적분}

$u = t,\ du = dt,\ dv = e^{-st}\,dt,\ v = -\frac{1}{s}e^{-st}$

\begin{equation}
\int_0^\infty t\, e^{-st}\, dt = \frac{1}{s}\int_0^\infty e^{-st}\, dt = \frac{1}{s^2}
\end{equation}

\subsubsection{종합}

\begin{equation}
\mathcal{L}[t^2] = \frac{2}{s} \cdot \frac{1}{s^2} = \frac{2}{s^3}
\end{equation}

\subsection{일반 공식}

같은 패턴으로 일반화하면:

\begin{empheq}[box=\fbox]{equation}
\mathcal{L}[t^n] = \frac{n!}{s^{n+1}}
\end{empheq}

\subsection{여기서 던질 만한 다섯 가지 통찰력 있는 질문}

\subsubsection{질문 1: 왜 $n!$이 등장하는가}

$n!$은 사실 \textbf{적분의 자연스러운 산물}이다. 다음 식은 \textbf{감마 함수}의 정의이다.

\begin{equation}
n! = \Gamma(n+1) = \int_0^\infty t^n e^{-t}\, dt
\end{equation}

즉 $\mathcal{L}[t^n]$ 계산은 사실상 $s=1$일 때 \emph{팩토리얼의 정의 자체}이다.

더 놀라운 것은 감마 함수 덕분에 분수 차수도 계산된다.

\begin{equation}
\mathcal{L}[t^{1/2}] = \frac{\Gamma(3/2)}{s^{3/2}} = \frac{\sqrt{\pi}}{2s^{3/2}}
\end{equation}

조합론(팩토리얼)과 미적분(감마 함수)과 원주율($\pi$)이 한 식에서 만나는 수학의 아름다운 다리 중 하나이다.

\subsubsection{질문 2: $s$ 미분과의 대칭성}

\begin{equation}
\frac{d}{ds}\left(\frac{1}{s}\right) = -\frac{1}{s^2}, \quad \frac{d^2}{ds^2}\left(\frac{1}{s}\right) = \frac{2}{s^3}, \quad \ldots
\end{equation}

이 패턴은 일반적인 쌍대성으로 정리된다.

\begin{equation}
\mathcal{L}[t \cdot f(t)] = -\frac{d}{ds}F(s)
\end{equation}

\textbf{원래 세계에서 $t$를 곱하는 것 = 변환된 세계에서 $s$로 미분하는 것}

\begin{center}
\begin{tabular}{ll}
\toprule
원래 세계 & 변환된 세계 \\
\midrule
미분 $f'(t)$ & 곱셈 $sF(s)$ \\
곱셈 $tf(t)$ & 미분 $-F'(s)$ \\
\bottomrule
\end{tabular}
\end{center}

완벽한 대칭. 이는 푸리에 변환의 \textbf{하이젠베르크 불확정성 원리}와도 연결되는 깊은 구조이다.

\subsubsection{질문 3: 역변환은 어떻게}

엄밀한 역변환 공식은 \textbf{브롬위치 적분}이다.

\begin{equation}
f(t) = \frac{1}{2\pi i}\int_{\gamma - i\infty}^{\gamma + i\infty} F(s)\, e^{st}\, ds
\end{equation}

복소평면의 수직선 위를 지나가는 적분이다. 라플라스 변환의 진정한 무대는 \textbf{실수가 아니라 복소수}임을 보여준다.

실용적으로는 거의 모든 경우에 \textbf{부분분수 분해}로 $F(s)$를 변환표의 조각들로 쪼개서 처리한다. 예:

\begin{equation}
\frac{1}{s^2 - 1} = \frac{1/2}{s-1} - \frac{1/2}{s+1} \xrightarrow{\mathcal{L}^{-1}} \frac{1}{2}e^t - \frac{1}{2}e^{-t} = \sinh t
\end{equation}

\subsubsection{질문 4: 수렴영역(ROC)}

$\mathcal{L}[e^t]$를 테일러 급수로 항별 변환하면 등비급수가 나오고, 수렴하려면 $s > 1$이어야 한다.

\begin{equation}
\mathcal{L}[e^{at}] = \frac{1}{s-a} \quad (\text{ROC: } s > a)
\end{equation}

라플라스 변환은 모든 $s$에서 정의되지 않는다 --- 특정 영역에서만 정의된다.

\begin{itemize}
\item $\mathcal{L}[1] = 1/s$ (ROC: $s > 0$)
\item $\mathcal{L}[e^{at}] = 1/(s-a)$ (ROC: $s > a$)
\item $\mathcal{L}[e^{t^2}]$: 어디서도 정의되지 않음 ($e^{t^2}$이 $e^{-st}$보다 빠르게 발산)
\end{itemize}

이 관점은 신호처리와 제어이론에서 \textbf{시스템의 안정성} 판별에 직접 쓰인다 --- 폴(pole)이 어디에 있느냐가 시스템이 안정한가를 결정한다.

\subsubsection{질문 5: 디랙 델타의 변환}

$\delta$의 체질(sifting property): $\int f(t)\delta(t)\,dt = f(0)$.

\begin{equation}
\mathcal{L}[\delta(t)] = \int_0^\infty \delta(t)\, e^{-st}\, dt = e^{0} = 1
\end{equation}

\textbf{모든 $s$에서 동일한 값} --- 즉 모든 지수함수 성분을 똑같이 포함한다는 의미. 시간 영역에서 한 점에 무한히 집중된 신호가, 변환 영역에서는 완벽하게 펴진 평탄한 분포가 된다. \textbf{시간-주파수 쌍대성의 가장 극단적 표현}.

이는 제어공학에서 \textbf{시스템의 임펄스 응답} 개념의 수학적 토대이다. 시스템에 순간적인 충격($\delta$)을 가했을 때 어떻게 반응하는가 --- 그 응답이 시스템의 모든 것을 알려준다.

\subsection{좋은 질문의 공통 구조}

\begin{itemize}
\item \textbf{패턴 인식}: ``왜 하필 이 모양이지?'' (질문 1, 4)
\item \textbf{대칭성}: ``한 방향이 됐는데, 반대 방향은?'' (질문 2, 3)
\item \textbf{경계 사례}: ``이 도구의 한계는 어디?'' (질문 5)
\end{itemize}

수학에서 가장 좋은 질문은 보통 이 세 갈래에서 나온다.

% ============================================================
\section{PDE 풀이의 세 갈래}
% ============================================================

\subsection{해석적 vs 수치적 vs ML 기반}

\begin{center}
\begin{tabular}{p{4cm}p{10cm}}
\toprule
방법 & 특성 \\
\midrule
\textbf{해석적} (적분변환, 좌표변환 등) & 답이 수식, 정확함. 단순한 PDE만 가능. \\
\textbf{수치적} (FDM, FVM, FEM) & 거의 모든 PDE 가능. 답이 격자 위 숫자. 차원의 저주 → 슈퍼컴퓨터 필요. \\
\textbf{ML 기반} (PINN, FNO 등) & 신경망 가중치 학습으로 해 함수 근사. 고차원ㆍ역문제에 강함. \\
\bottomrule
\end{tabular}
\end{center}

\subsection{수치적 방법의 종류}

\begin{itemize}
\item \textbf{FDM (Finite Difference Method)}: $\frac{\partial u}{\partial x} \approx \frac{u_{i+1} - u_i}{\Delta x}$ --- 미분을 차분으로 근사
\item \textbf{FVM (Finite Volume Method)}: 영역을 작은 체적으로 쪼개고, 각 체적의 경계를 통과하는 흐름의 보존 법칙을 적용
\item \textbf{FEM (Finite Element Method)}: 영역을 작은 조각(element)으로 쪼개고, 각 조각 안에서 간단한 함수(보통 다항식)로 해를 근사
\end{itemize}

\subsection{차원의 저주}

3차원 공간 + 시간의 PDE에서 공간을 $N$개씩 쪼개면:

\begin{itemize}
\item 격자점 수: $N^3$
\item 시간 스텝 포함: 최소 $N^4$
\item $N = 1000$일 때: $10^{12}$번의 기본 계산 (그것도 한 스텝당)
\item 비선형이면 매 스텝마다 거대한 행렬 풀이 추가
\end{itemize}

이것이 슈퍼컴퓨터(일본 \emph{후가쿠}, 미국 \emph{프론티어}, \emph{El Capitan})가 필요한 본질이다.

\subsection{PINN: Physics-Informed Neural Networks}

2017--2019년경 라이시(Maziar Raissi), 카르니아다키스(George Karniadakis) 등이 제안.

\subsubsection{핵심 아이디어}

해 함수 $u(x, t)$ 자체를 신경망으로 모델링한다.

\begin{equation}
u(x, t) \approx u_\theta(x, t) \quad (\theta: \text{신경망 가중치})
\end{equation}

손실함수:

\begin{equation}
\mathcal{L}(\theta) = \mathcal{L}_{\text{PDE}} + \mathcal{L}_{\text{BC}} + \mathcal{L}_{\text{IC}}
\end{equation}

열방정식 $u_t = u_{xx}$의 경우:

\begin{equation}
\mathcal{L}_{\text{PDE}} = \frac{1}{N}\sum_{i=1}^N \left|\frac{\partial u_\theta}{\partial t}(x_i, t_i) - \frac{\partial^2 u_\theta}{\partial x^2}(x_i, t_i)\right|^2
\end{equation}

\subsubsection{핵심 기술: 자동미분(Automatic Differentiation)}

신경망의 미분을 수치미분이 아닌 \textbf{정확한 미분}으로 계산할 수 있다. 신경망은 합성함수 덩어리라서 연쇄법칙을 자동으로 적용하여 미분을 정확히 계산한다. PyTorch, TensorFlow의 핵심 기능.

\subsubsection{PINN vs 전통 수치해석}

\begin{center}
\begin{tabular}{lll}
\toprule
& 전통 FVM/FEM & PINN \\
\midrule
미지수 & 격자점의 값 & 신경망 가중치 \\
해의 표현 & 격자 위 숫자들 & 연속 함수 \\
격자 필요 & 필수 & 불필요 \\
차원의 저주 & 심각 & 완화 \\
역문제 & 어려움 & 자연스러움 \\
정확도 보장 & 격자 정밀화로 가능 & 이론적 보장 부족 \\
\bottomrule
\end{tabular}
\end{center}

\subsubsection{PINN의 강점과 한계}

\textbf{강점}:
\begin{itemize}
\item 고차원 PDE 가능 (10차원, 100차원도). 금융수학의 옵션 가격 결정 등.
\item 역문제: ``관측 데이터로 PDE의 계수를 찾아라'' 같은 문제에 자연스러움.
\end{itemize}

\textbf{한계}:
\begin{itemize}
\item 학습 불안정. 다중 손실항의 균형 어려움.
\item 정확도 보장 부족.
\item 복잡한 형상ㆍ다물리 문제에는 전통 방법이 우세.
\item 매번 새 PDE마다 처음부터 재학습.
\end{itemize}

% ============================================================
\section{미분과 차분, 특성곡선법, 그린 함수}
% ============================================================

\subsection{미분과 차분의 차이}

\subsubsection{차분 (Difference)}

\begin{equation}
\frac{\Delta f}{\Delta x} = \frac{f(x + h) - f(x)}{h}
\end{equation}

\textbf{$h$가 유한한 값.} 두 점에서의 함숫값 차이를 거리로 나눈 것. 평균변화율.

\subsubsection{미분 (Differentiation)}

\begin{equation}
\frac{df}{dx} = \lim_{h \to 0} \frac{f(x + h) - f(x)}{h}
\end{equation}

차분에서 $h$를 \textbf{0으로 보낸 극한}. 순간변화율, 접선 기울기.

\subsubsection{비유}

\begin{itemize}
\item \textbf{차분(평균 속도)}: ``10초 동안 200m를 달렸으니 평균 시속 72km''
\item \textbf{미분(순간 속도)}: 자동차 계기판이 보여주는 값
\end{itemize}

물리적으로 순간 속도는 측정 불가능하다. 모든 측정은 짧은 시간 간격을 필요로 한다.

\subsubsection{컴퓨터의 한계}

컴퓨터는 극한을 직접 다룰 수 없다. PDE의 연속적인 미분을 이산적인 차분으로 \textbf{바꿔치기}해서 푸는 것이 수치해석의 출발이다.

\begin{equation}
\frac{\partial u}{\partial x}\bigg|_{x_i} \approx \frac{u_{i+1} - u_i}{\Delta x}
\end{equation}

이때 발생하는 오차가 \textbf{이산화 오차(discretization error)}.

\begin{historybox}[역사적 사실]
미분보다 차분이 먼저였다. 17세기 수학자들(윌리스, 페르마)이 유한차분으로 곡선의 기울기와 면적을 다루고 있었다. 뉴턴/라이프니츠의 혁명은 이 유한차분에 \emph{극한}을 도입하여 순간적 양을 정의한 것이다. 21세기 컴퓨터가 다시 차분으로 돌아간 것은 수학사의 원점 회귀이기도 하다.
\end{historybox}

\subsection{특성곡선법 (Method of Characteristics)}

\subsubsection{핵심 아이디어}

PDE의 \textbf{특별한 곡선들}을 찾아, 그 곡선을 따라가면 PDE가 단순한 ODE로 변하는 방법. 이 곡선을 \textbf{특성곡선(characteristics)}이라 한다.

비유: 험한 산에서 길을 찾을 때 물이 흐르는 길을 따라가면 자연스럽게 아래로 내려간다. 특성곡선은 PDE의 ``물길''이다.

\subsubsection{가장 단순한 예: 이송방정식}

\begin{equation}
\frac{\partial u}{\partial t} + c \frac{\partial u}{\partial x} = 0
\end{equation}

특성곡선 $x = ct + x_0$ 위에서:

\begin{equation}
\frac{du}{dt}\bigg|_{\text{특성선 위}} = \frac{\partial u}{\partial t} + \frac{\partial u}{\partial x}\cdot c = 0
\end{equation}

즉 특성선 위에서 $u$는 상수. 따라서 초기 모양이 속도 $c$로 평행이동:

\begin{equation}
u(x, t) = u_0(x - ct)
\end{equation}

\subsubsection{파동방정식과의 연결}

\begin{equation}
\left(\frac{\partial}{\partial t} - c\frac{\partial}{\partial x}\right)\left(\frac{\partial}{\partial t} + c\frac{\partial}{\partial x}\right)u = 0
\end{equation}

두 개의 1차 이송 연산자로 인수분해되며, 두 개의 특성족 $x \pm ct = \text{const}$에 해당한다. 해는 두 방향으로 흘러가는 신호의 합 $u(x, t) = f(x - ct) + g(x + ct)$ --- \textbf{달랑베르 해의 진짜 정체}이다.

\subsubsection{비선형 확장: 충격파}

비선형 이송방정식:

\begin{equation}
\frac{\partial u}{\partial t} + u\frac{\partial u}{\partial x} = 0
\end{equation}

특성선의 속도가 $u$ 자신에 의존하므로, 값이 큰 곳의 특성선이 빠르게 움직여 \textbf{서로 교차}할 수 있다. 자연은 이를 \textbf{충격파(shock wave)}로 해결한다 --- 소닉붐, 교통체증의 정체파가 모두 같은 수학적 메커니즘.

\subsection{그린 함수 (Green's Function)}

\subsubsection{역사}

\begin{historybox}[조지 그린 (George Green, 1793--1841)]
영국 노팅엄의 제분소 주인의 아들. 초등학교만 1년 다녔고 나머지는 모두 독학. 30대 후반에 자비로 출판한 \emph{전기와 자기의 수학적 분석에 관한 에세이(1828)}에서 그린 함수 개념을 처음 제시. 당시에는 거의 무시당했고, 켈빈 경(William Thomson)이 1840년대에 재발견하여 세상에 알렸다. 그린은 자신의 도구의 위력을 알아보지 못한 채 48세에 세상을 떴다.
\end{historybox}

\subsubsection{핵심 아이디어: 점원천의 응답}

\textbf{복잡한 입력을 점입력들로 분해할 수 있다면, 각 점입력의 응답을 알아두고 합치면 전체 응답을 얻는다.}

푸아송 방정식 예시:

\begin{equation}
-\nabla^2 \phi = \rho
\end{equation}

먼저 가장 단순한 원천(점원천)에 대한 응답 $G$를 구한다.

\begin{equation}
-\nabla^2 G(x, x') = \delta(x - x')
\end{equation}

이 $G(x, x')$가 \textbf{그린 함수}이다. 3차원에서:

\begin{equation}
G(x, x') = \frac{1}{4\pi|x - x'|} \quad (\text{쿨롱 전위})
\end{equation}

\subsubsection{중첩 원리의 마법}

임의의 전하 분포는 점전하들의 합으로 분해된다.

\begin{equation}
\rho(x) = \int \rho(x')\delta(x - x')\, dx'
\end{equation}

방정식이 선형이므로 응답을 중첩할 수 있다.

\begin{empheq}[box=\fbox]{equation}
\phi(x) = \int G(x, x')\,\rho(x')\, dx'
\end{empheq}

\textbf{한 번 $G$를 구해두면, 어떤 원천이 와도 적분 한 번으로 끝.}

\subsubsection{비유: 콘서트홀의 음향}

콘서트홀의 어느 자리에서 한 번 박수를 친다. 청중석에서 어떻게 들리는지 측정한 것이 \textbf{임펄스 응답} = 그린 함수의 물리적 형제. 이걸 알아두면 그 자리에서 어떤 음악을 연주해도 청중석의 소리를 계산할 수 있다.

이것이 신호처리에서 \textbf{시스템의 임펄스 응답 = 시스템의 완전한 정체성}이라는 원리이며, $\mathcal{L}[\delta] = 1$과 직접 연결된다.

\subsubsection{그린 함수의 응용}

\begin{itemize}
\item \textbf{전자기학}: 점전하의 쿨롱 포텐셜
\item \textbf{양자장론}: 입자의 \emph{전파인자(propagator)}. 파인만 다이어그램의 각 선이 그린 함수.
\item \textbf{열전달}: 한 점의 순간 열 펄스의 응답
\item \textbf{확률론과 금융수학}: 전이확률, 옵션 가격 결정
\item \textbf{의료영상}: CT, MRI 재구성
\end{itemize}

\subsection{세 도구의 통합 관점}

\begin{center}
\begin{tabular}{lll}
\toprule
도구 & 철학 & 시대 정신 \\
\midrule
특성곡선법 & PDE의 물리적 흐름을 따라가자 & 18--19세기, 물리적 직관 \\
그린 함수 & 점원천 응답을 알면 모두 안다 & 19세기, 선형성과 중첩 \\
차분 (수치해석) & 극한을 포기하고 격자에서 풀자 & 20세기, 컴퓨터 \\
\bottomrule
\end{tabular}
\end{center}

% ============================================================
\section{FNO: 푸리에 신경연산자}
% ============================================================

\subsection{가장 큰 도약: 함수가 아니라 연산자를 학습하라}

\subsubsection{PINN vs FNO}

\begin{itemize}
\item \textbf{PINN}: 특정 PDE의 \emph{특정 해 함수} $u(x, t)$를 신경망으로 표현. 초기조건이 바뀌면 처음부터 재학습.
\item \textbf{FNO}: PDE의 \emph{연산자(operator) 자체}를 학습. 임의의 입력 함수에 대한 해 함수를 출력하는 \emph{규칙}을 배운다.
\end{itemize}

\begin{equation}
\mathcal{G}: \underbrace{a(x)}_{\text{입력 함수}} \mapsto \underbrace{u(x)}_{\text{해 함수}}
\end{equation}

\textbf{비유}: PINN은 ``피아노 한 곡을 외워 치는 학생'', FNO는 ``악보 읽는 법을 배운 피아니스트''. 한 번 학습되면 새 입력에도 즉시 ``연주''할 수 있다.

\subsubsection{연산자 학습의 역사}

\begin{itemize}
\item \textbf{2019}: DeepONet (Lu et al.) --- 첫 신경연산자
\item \textbf{2020}: FNO (Li et al.) --- 푸리에 공간에서 적분 커널을 매개변수화
\end{itemize}

\subsection{이론적 토대: 그린 함수의 부활}

선형 PDE의 해는 그린 함수와 입력의 합성곱:

\begin{equation}
u(x) = \int G(x, x')\, a(x')\, dx'
\end{equation}

일반화하면 임의의 연산자를 적분 커널로 표현 가능.

\begin{equation}
(\mathcal{G}a)(x) = \int \kappa(x, x')\, a(x')\, dx'
\end{equation}

\begin{insightbox}[FNO의 핵심 발상]
\textbf{이 커널 $\kappa$를 신경망으로 학습하자.}

즉 FNO는 ``\textbf{학습 가능한 그린 함수}''이다. 어떤 PDE인지 명시적으로 모르더라도, 데이터를 통해 그 PDE의 그린 함수에 해당하는 커널을 신경망이 찾아낸다. 19세기 그린 함수의 21세기 부활.
\end{insightbox}

\subsection{왜 푸리에인가: 합성곱 정리}

직접 적분 $\int \kappa(x, x')a(x')dx'$를 계산하면 격자점 $N$개에 대해 $O(N^2)$. 3차원이면 $O(N^6)$ --- 차원의 저주 부활.

\textbf{합성곱 정리}:

\begin{equation}
\widehat{f * g} = \hat{f} \cdot \hat{g}
\end{equation}

실공간에서의 합성곱 = 푸리에 공간에서의 단순 곱셈. FFT는 $O(N\log N)$.

\begin{empheq}[box=\fbox]{equation}
\underbrace{O(N^2)}_{\text{실공간 합성곱}} \;\Longrightarrow\; \underbrace{O(N\log N)}_{\text{푸리에 경유 합성곱}}
\end{empheq}

\textbf{그린 함수와의 적분을 푸리에 공간에서의 단순 곱셈으로 바꿔치기.}

\subsection{FNO 아키텍처}

FNO는 크게 세 부분으로 구성된다.

\begin{enumerate}
\item \textbf{리프팅(Lifting) $P$}: 입력 함수를 고차원 채널 공간으로 끌어올림.
\item \textbf{푸리에 층 (Fourier Layer)} (반복, 보통 4--8회): 핵심부.
\item \textbf{투영(Projection) $Q$}: 다시 출력 차원으로 압축.
\end{enumerate}

각 푸리에 층의 네 단계:

\begin{enumerate}
\item \textbf{FFT}: 푸리에 공간으로 변환
\item \textbf{저차 모드 곱셈}: 낮은 주파수 모드에만 학습 가능한 복소수 가중치 $R$ 곱셈 (\textbf{학습의 핵심})
\item \textbf{역 FFT}: 실공간으로 복귀
\item \textbf{잔차 결합}: 선형 변환 $W$를 옆길로 더함 (ResNet 스타일 skip connection)
\end{enumerate}

\subsubsection{저차 모드 절단의 의미}

학습 가능한 푸리에 필터를 차원당 가장 낮은 $m$개 모드로 제한. 자연 현상이 대체로 부드럽기 때문에(저주파 성분이 지배적이기 때문에) 잘 작동한다. 부드러움 정규화 역할도 한다.

\subsection{진짜 마법: 해상도 무관성}

\begin{insightbox}[Discretization Invariance]
$64 \times 64$ 격자로 학습한 신경망이 즉시 $1024 \times 1024$ 격자에서 잘 작동한다.

\textbf{푸리에 모드는 격자 크기와 무관한 함수}이기 때문이다. 모드 5번은 격자가 64개든 1024개든 같은 사인파다. 따라서 저차 모드 위의 가중치 $R$은 격자 해상도와 분리되어 있다.

전통 시뮬레이션은 해상도가 두 배 되면 계산이 16배 폭증하는 반면, FNO는 학습 비용에 거의 변화가 없다.
\end{insightbox}

\subsection{실전: FourCastNet과 NVIDIA Earth-2}

\subsubsection{FourCastNet (2022)}

NVIDIA, 로렌스 버클리 국립연구소, 미시간 대학, 라이스 대학 공동 개발.

\begin{itemize}
\item 해상도: $0.25^\circ$ (적도 부근 $30\text{km} \times 30\text{km}$, $720 \times 1440$ 픽셀)
\item 속도: 전통 NWP 대비 \textbf{45,000배 빠름}
\item 정확도: ECMWF IFS 시스템과 비견될 수준 (IFS와 직접 비교 가능한 최초의 AI 모델)
\item 일주일치 예보 생성: \textbf{2초}
\end{itemize}

\subsubsection{FourCastNet 3 (2025--2026)}

\begin{itemize}
\item 단일 NVIDIA H100 GPU에서 60일치 0.25도, 6시간 해상도 예보 = \textbf{4분}
\item GenCast 대비 8배, IFS-ENS 대비 60배 빠름
\item 60일 확장 리드 타임에서도 스펙트럼 특성 보존
\item 운영 환경 투입: DTN의 사이클론 추적 모델 (2025년 6월), Eni의 가스 수요 예측, GCL의 광전지 예측 시스템 등
\end{itemize}

\subsection{한계}

\begin{itemize}
\item \textbf{주기적 경계 가정}: 푸리에 변환은 본질적으로 주기 함수 가정. 복잡한 형상에는 어려움. (해결책: Spherical FNO, Geo-FNO, Domain Decomposition FNO)
\item \textbf{저주파 편향}: 전역 푸리에 필터와 절단 때문에 고주파 현상(날카로운 경계, 충격파)을 놓치기 쉬움.
\item \textbf{모드 절단의 트레이드오프}: 편향-분산. 적절한 모드 수 선택 자체가 학습 작업.
\end{itemize}

그래도 2025--2026년 현재 신경연산자 분야의 \textbf{표준 아키텍처}이다.

% ============================================================
\section{결론: 280년의 흐름}
% ============================================================

\subsection{한 줄로 꿴 여정}

\begin{center}
\begin{tabular}{p{2.5cm}p{11.5cm}}
\toprule
연도 & 사건 \\
\midrule
\textbf{1747} & \textbf{달랑베르} --- 좌표변환으로 파동방정식 풀이. 특성곡선의 첫 등장. \\
\textbf{1782--1809} & \textbf{라플라스} --- 함수를 지수함수들의 합으로 분해. 미분을 곱셈으로. \\
\textbf{1807--1822} & \textbf{푸리에} --- 함수를 사인파의 합으로 분해. \\
\textbf{1828} & \textbf{그린} --- 점원천 응답으로 모든 응답을 안다. 적분 커널의 발견. \\
\textbf{1880년대} & \textbf{헤비사이드} --- 연산자 미적분, 전기회로 응용. \\
\textbf{20세기 중반} & \textbf{컴퓨터 등장} --- 미분을 차분으로. FDM, FVM, FEM. 슈퍼컴퓨터의 시대. \\
\textbf{2017} & \textbf{PINN} --- 자동미분으로 PDE를 손실함수의 제약으로. 신경망이 해 함수가 된다. \\
\textbf{2020} & \textbf{FNO} --- 연산자 자체를 학습. 푸리에 공간에서 그린 함수에 해당하는 적분 커널을 데이터로부터 발견. \\
\textbf{2025--2026} & \textbf{FourCastNet 3} --- 운영 환경 투입. 일기예보가 단일 GPU에서 분 단위로. \\
\bottomrule
\end{tabular}
\end{center}

\subsection{불변의 사상}

도구는 바뀌었어도 \textbf{바꿔치기의 정신}은 변하지 않았다.

\begin{itemize}
\item 미분을 \textbf{차분}으로 바꾸고 (수치해석)
\item 미분을 \textbf{곱셈}으로 바꾸고 (적분변환)
\item PDE를 \textbf{특성곡선 위의 ODE}로 바꾸고 (특성곡선법)
\item PDE를 \textbf{그린 함수와의 합성곱}으로 바꾸고 (그린 함수)
\item PDE를 \textbf{신경망의 자동미분과 손실함수}로 바꾸고 (PINN)
\item PDE를 \textbf{학습된 푸리에 공간의 곱셈}으로 바꾼다 (FNO)
\end{itemize}

280년의 수리물리학과 계산과학의 흐름은 결국 단 하나의 사상으로 수렴한다:

\begin{center}
\fbox{\parbox{0.8\textwidth}{\centering \large
\textbf{``다루기 어려운 미분 연산을\\다루기 쉬운 무언가로 바꿔치기하는 기술의 발전사''}
}}
\end{center}

그리고 그 마지막 정거장에서, 푸리에가 처음 던진 ``함수를 진동의 합으로 분해한다''는 발상은 21세기 일기예보의 심장부에서 다시 살아 움직이고 있다. 200년 전의 수학이 살아남기는커녕 더 깊이 작동하고 있다 --- 이것이 수학의 진짜 아름다움이다.

% ============================================================
\appendix
\section{더 깊이 들어가고 싶다면}
% ============================================================

이 노트에서 다룬 주제를 더 파고들 수 있는 방향들:

\begin{itemize}
\item \textbf{DeepONet}: FNO의 사촌. 다른 접근(branch + trunk 네트워크)으로 연산자 학습.
\item \textbf{Diffusion 기반 PDE 풀이}: 2024--2025년 가장 활발한 흐름. 확산 모델을 PDE 해결에 응용.
\item \textbf{신경연산자의 이론적 보장}: Operator Universal Approximation Theorem.
\item \textbf{Spherical FNO, Geo-FNO}: 구면 또는 불규칙 도메인에서의 FNO.
\item \textbf{Differentiable Physics}: 시뮬레이션 전체를 미분 가능하게 만들어 역설계.
\item \textbf{NVIDIA PhysicsNeMo, Earth-2}: Physics-ML의 산업적 구현.
\item \textbf{함수해석학과 분포 이론}: 슈바르츠, 힐베르트 공간, 소볼레프 공간.
\item \textbf{복소해석학과 라플라스 변환의 역변환}: 브롬위치 적분의 본격 다루기.
\end{itemize}

\section{컴파일 안내}

이 문서는 XeLaTeX로 컴파일해야 한다 (한국어 처리를 위해 \texttt{kotex} 패키지 사용).

\begin{verbatim}
xelatex pde_notes.tex
xelatex pde_notes.tex  # 목차 갱신을 위해 두 번 실행
\end{verbatim}

필요한 시스템:
\begin{itemize}
\item TeX Live 2020 이상 (또는 MiKTeX)
\item 한국어 폰트 (Nanum, Noto Sans CJK KR 등)
\item 패키지: \texttt{kotex}, \texttt{amsmath}, \texttt{amssymb}, \texttt{geometry}, \texttt{hyperref}, \texttt{xcolor}, \texttt{tcolorbox}, \texttt{booktabs}, \texttt{enumitem}, \texttt{parskip}
\end{itemize}

\end{document}
